\section{Methods}

\subsection{Datasets}

We reassembled and reprocessed six previously published comparison studies, each of
which profiled a common biological sample with both the Parse Biosciences Evercode
Whole Transcriptome (WT) split-pool platform and the 10x Genomics Chromium Single
Cell 3$'$ droplet platform (Table~\ref{tab:datasets}). Raw reads were obtained from
the NCBI Sequence Read Archive (SRA) or the European Nucleotide Archive (ENA); the
accession for every assay is recorded in the accompanying code repository. The panel
comprised mouse thymocytes (Analysis~2;
Evercode WT v2 and WT Mini v2 vs.\ Chromium Next GEM v3), human PBMCs from two donors
(Analysis~3; Evercode WT v2 vs.\ Chromium Next GEM v3), a human/mouse cell-line
``barnyard'' mixture of K562 and mESC cells (Analysis~4; Evercode WT Mini v2 vs.\
Chromium Next GEM v3), frozen human PBMCs (Analysis~5; Evercode WT v2 vs.\ Chromium
Next GEM v3), a multiplexed human glucocorticoid-resistant E/R+ ALL Reh cell line
under six drug perturbations (Analysis~6; Evercode WT Mini v1 vs.\ Chromium Next GEM
v3), and resting/activated frozen human PBMCs (Analysis~7; Evercode WT Mini v3 vs.\
Chromium v4). To remove sequencing depth as a confounder, all assays within an
analysis were processed through a single unified pipeline and compared at matched
read depth (see \textit{Read-depth matching}).

All processing was orchestrated by a custom Python package (\texttt{XvP\_utils})
built on a conda environment; key tool versions are given in Table~\ref{tab:software}.

\begin{table*}[t]
\centering
\footnotesize
\setlength{\tabcolsep}{5pt}
\caption{Datasets. Each study profiled a common sample with both the Parse Evercode
and 10x Chromium platforms. Sequencing accessions for every assay are provided in the
accompanying code repository. Analysis~7 is an in-house, currently unpublished
dataset.}
\label{tab:datasets}
\begin{tabular}{@{}p{0.04\textwidth} p{0.33\textwidth} p{0.26\textwidth} p{0.26\textwidth}@{}}
\hline
\# & Sample & Parse Evercode & 10x Chromium \\
\hline
2 & Mouse thymocytes, two female mice & WT v2; WT Mini v2 & Next GEM v3 (+ TotalSeq-B hashtags) \\
3 & Human PBMCs, two healthy donors & WT v2 & Next GEM v3 \\
4 & Human K562 / mouse mESC barnyard mixture & WT Mini v2 & Next GEM v3 \\
5 & Frozen human PBMCs & WT v2 & Next GEM v3 \\
6 & Human E/R+ ALL Reh cell line, six drug perturbations & WT Mini v1 & Next GEM v3 (+ MULTI-seq) \\
7 & Resting/activated frozen human PBMCs & WT Mini v3 & Chromium v4 (+ MULTI-seq) \\
\hline
\end{tabular}
\end{table*}

\subsection{Read retrieval and library assembly}

SRA runs were downloaded with \texttt{prefetch} and converted to FASTQ with
\texttt{fasterq-dump} (\texttt{-{}-split-files -{}-skip-technical}; SRA-Tools
v3.0.3~\cite{leinonen2011sra}) and compressed with \texttt{pigz}. ENA runs were
streamed directly from the ENA FTP server as gzipped FASTQ. For assays distributed as
multiple sublibraries, all sublibraries belonging to a single assay were combined
into one paired-end FASTQ with \texttt{splitcode -{}-remultiplex} (splitcode
v0.30.0~\cite{sullivan2024splitcode}), preserving the original
read-1/read-2 structure. Because a Parse sublibrary is resolved by its sequencing
index rather than by the combinatorial barcodes themselves, identical round-1/2/3
barcode combinations occurring in different sublibraries label different cells. Parse
sublibraries were therefore combined with \texttt{-{}-bclen=4}, which prepends a
distinct 4\,bp sequence to read~2 of every sublibrary; runs originating from the same
sublibrary were assigned the same sequence. This synthetic barcode was carried through
the remainder of the pipeline as a fourth cell barcode.

\subsection{Reference genomes and indices}

Reference genome FASTA and GTF annotations were obtained from Ensembl~\cite{ensembl}. Human assays
used GRCh38 (release~115) and the mouse thymocyte assay used GRCm39 (release~115).
For the barnyard analysis a combined reference was constructed by concatenating the
human (GRCh38, release~112) and mouse (GRCm39, release~115) genomes; every chromosome
name, \texttt{gene\_id}, and \texttt{gene\_name} was prefixed with a species tag
(\texttt{HUMAN\_}~/~\texttt{MOUSE\_}) so that reads could be assigned unambiguously to
a species of origin.

For pseudoalignment, a nascent-and-mature (nac) index was built for each reference
with \texttt{kb ref -{}-workflow nac} (kb-python v0.29.3~\cite{sullivan2023kb},
wrapping kallisto v0.51.1~\cite{bray2016kallisto} and bustools
v0.44.1~\cite{melsted2021bustools}). The nac workflow indexes both spliced (mature)
and unspliced
(nascent) transcript sequences, allowing intronic reads to be quantified separately
from exonic reads.

\subsection{Parse split-pool barcode handling}

Parse Evercode encodes cell identity through three rounds of combinatorial barcoding
carried on read~2, with the round-1 (RT) barcode existing in two chemistries within
each well: an anchored poly(dT) primer (polyT) and a random-hexamer primer (randO).
Because the polyT and randO primers in a given well label the same cell, they must
be reconciled during quantification.

Barcode whitelists and the kallisto technology string for each kit/chemistry were
generated programmatically from the Parse barcode reference
(\texttt{Configs/parse\_info/}), keyed on the kit name (e.g.\ \texttt{WT\_v2},
\texttt{WT\_mini\_v3}) and, where applicable, the subset of sample wells used. All
read-2 barcode and UMI coordinates were offset by 4\,bp to account for the sublibrary
barcode prepended during library assembly, and the sublibrary barcode itself was
registered as an additional cell barcode occupying the first 4\,bp of read~2. Reads
were then processed with \texttt{splitcode} (v0.30.0) in two steps: (i)~reads were
filtered to those carrying an expected round-1 barcode at the kit-specific position,
allowing one mismatch (Hamming distance $\leq 1$); and (ii)~the retained reads were
partitioned into polyT and randO subsets for separate downstream analysis.

Filtered Parse reads were quantified with \texttt{kb count -{}-workflow nac} using
the kit-specific technology string, forward strandedness
(\texttt{-{}-strand=forward}), and single-end parity
(\texttt{-{}-parity=single}). A per-round barcode on-list (\texttt{-w onlist.txt})
was supplied for barcode error correction, listing the sublibrary barcode alongside
the three combinatorial rounds so that cells were resolved by sublibrary as well as
by barcode combination. A replacement map
(\texttt{-r replace.txt}) collapsed each randO round-1 barcode onto its polyT
counterpart in the same well so that both primer chemistries contributed counts to a
single cell. Three count matrices were produced per Parse assay: the full filtered
set, the polyT-only subset, and the randO-only subset.

\subsection{10x Genomics quantification}

10x reads were quantified with \texttt{kb count -{}-workflow nac} using the
appropriate technology string (\texttt{10XV3} for Chromium Next GEM v3 assays,
\texttt{10XV4} for the Chromium v4 assay) and the matching cell-barcode on-list. As
with the Parse data, the nac workflow yielded mature, nascent, and ambiguous count
layers.

For the multiplexed datasets (Analyses~6 and 7), the 10x hashtag/MULTI-seq oligo
libraries were quantified separately with the kb-python \texttt{kite} (feature
barcoding) workflow; read~2 was trimmed to the tag length with \texttt{seqtk
trimfq}~\cite{seqtk} prior to counting.

\subsection{Read-depth matching}

To compare technologies at equal sequencing effort, processed read counts were
tabulated for every assay in an analysis (the full filtered FASTQ for each 10x and
Parse assay, plus the Parse polyT and randO subsets) and the minimum across assays
was taken as the target depth. Each FASTQ was then downsampled to this common depth
with \texttt{seqtk sample} using a fixed random seed (42) and re-quantified with
kb-python as above. The 10x, full-Parse, polyT, and randO read sets were each
subsampled to the same depth so that all per-gene and per-cell comparisons reflect
differences in chemistry rather than in coverage.

\subsection{Cell calling and quality control}

Downstream analysis was performed in Python with Scanpy
(v1.11.4~\cite{wolf2018scanpy}) and AnnData (v0.12.2~\cite{virshup2024anndata}). Each
kb-python h5ad matrix was loaded, Ensembl gene IDs were mapped to
gene symbols, and genes/cells with zero counts were removed. Per-cell and per-gene
summary statistics were computed, including total UMI counts, genes detected, and
the fraction of counts attributable to nascent (unspliced), mature (spliced), and
ambiguous molecules. Alignment statistics were drawn from kb-python's
\texttt{run\_info.json} and, where STAR alignments were generated, from STARsolo's
\texttt{Log.final.out}.

Cells were called from each assay's barcode-rank (knee) plot: the UMI distribution
was ranked and an assay-specific minimum-UMI threshold was applied to separate cells
from empty droplets/barcodes, retaining all barcodes above the threshold.

Gene-level annotations were assembled once per reference. Gene biotype (protein
coding, lncRNA, mitochondrial, ribosomal, and OXPHOS classes, identified by biotype
and gene-symbol prefixes) was parsed from the Ensembl GTF, while transcript length
and GC content were computed directly from the kallisto cDNA FASTA and averaged
across isoforms per gene---guaranteeing complete coverage for every gene in the
index. Per-cell biotype fractions (percent protein-coding, mitochondrial,
ribosomal, and lncRNA counts) were added from these annotations.

Library complexity / sequencing saturation was estimated for each assay by
extrapolating the per-gene count distribution with PreSeq
\texttt{lc\_extrap}~\cite{daley2013preseq}.

\subsection{Gene-body coverage}

To assess 5$'$/3$'$ coverage bias, depth-matched reads were aligned to the genome
with STARsolo (STAR v2.7.11b~\cite{dobin2013star}). 10x reads were aligned in
\texttt{CB\_UMI\_Simple}
mode with the 10x cell-barcode whitelist; Parse reads were aligned in
\texttt{CB\_UMI\_Complex} mode with the three per-round Parse whitelists and
barcode/UMI positions derived from the kit technology string (1-mismatch barcode
matching). Parse alignments were generated separately for the full, polyT, and randO
read subsets. Coordinate-sorted BAM files were indexed with SAMtools
(v1.21~\cite{danecek2021samtools}). A BED12 model was produced from the GTF with
\texttt{gffread}~\cite{pertea2020gffread} and subset to housekeeping transcripts from
the HRT Atlas~\cite{hounkpe2021hrt}; transcript-body coverage was then computed across
all assays with RSeQC \texttt{geneBody\_coverage.py} (RSeQC v5.0.4~\cite{wang2012rseqc}).

\subsection{Per-gene composition comparison}

For each technology pair, the per-gene share of total counts (gene counts as a
percentage of all counts in that assay) was compared. Gene \texttt{var} tables were
merged on shared gene identifiers (genes detected in only one assay were retained
with a value of zero), and an ordinary least-squares regression of one assay's
per-gene percentage on the other's was fit with statsmodels~\cite{seabold2010statsmodels}. Cook's distance was
computed for every gene to flag those contributing disproportionately to the
between-technology difference, and a Gaussian kernel density estimate was used to
colour points by local density. Outlier genes were summarised by biotype, gene
length, GC content, and gene-name prefix.

\subsection{Cross-study differential composition and functional enrichment}

To identify genes that were consistently over- or under-represented by one technology
across studies, a pseudobulk differential-expression analysis was performed with
edgePython, a Python implementation of the edgeR
workflow~\cite{robinson2010edger,chen2016edger}. Counts were summed across all cells
within each assay to form one pseudobulk profile per sample, and profiles were
restricted to genes detected in every sample. A DGEList was constructed and normalised
by the trimmed mean of M-values (TMM) method~\cite{robinson2010tmm}. A design matrix of the form
\texttt{$\sim$ tissue + technology} was used so that the technology effect was
estimated while controlling for biological source; lowly expressed genes were removed
with \texttt{filterByExpr}. Negative-binomial dispersions were estimated
(\texttt{estimateDisp}), a quasi-likelihood negative-binomial GLM was fit
(\texttt{glmQLFit}), and the technology coefficient was tested with the
quasi-likelihood F-test (\texttt{glmQLFTest}). P-values were adjusted by the
Benjamini--Hochberg procedure, and genes with FDR below the chosen threshold were
called as differentially represented between technologies.

Genes consistently enriched in each technology were then tested for functional
enrichment with gseapy's~\cite{fang2023gseapy} Enrichr~\cite{kuleshov2016enrichr} interface against the GO Biological Process, GO
Molecular Function, GO Cellular Component, KEGG, and MSigDB Hallmark gene-set
libraries, using all genes retained in the comparison as the statistical background.
Terms with a Benjamini--Hochberg adjusted p-value below 0.05 were reported.

\begin{table}[ht]
\centering
\caption{Software.}
\label{tab:software}
\begin{tabular}{lll}
\hline
Tool & Version & Use \\
\hline
SRA-Tools~\cite{leinonen2011sra} & 3.0.3 & SRA download / FASTQ conversion \\
pigz & 2.4 & FASTQ compression \\
splitcode~\cite{sullivan2024splitcode} & 0.30.0 & Library remultiplexing, Parse barcode filtering/splitting \\
kb-python~\cite{sullivan2023kb} & 0.29.3 & Pseudoalignment and UMI counting (nac workflow) \\
\quad kallisto / bustools~\cite{bray2016kallisto,melsted2021bustools} & 0.51.1 / 0.44.1 & (bundled with kb-python) \\
seqtk~\cite{seqtk} & 1.2-r94 & Read subsampling and R2 trimming \\
STAR / STARsolo~\cite{dobin2013star} & 2.7.11b & Genome alignment for gene-body coverage \\
SAMtools~\cite{danecek2021samtools} & 1.21 & BAM indexing \\
gffread~\cite{pertea2020gffread} & 0.12.7 & GTF $\rightarrow$ BED12 conversion \\
RSeQC~\cite{wang2012rseqc} & 5.0.4 & Gene-body coverage \\
PreSeq~\cite{daley2013preseq} & 3.2.0 & Library-complexity extrapolation \\
Scanpy / AnnData~\cite{wolf2018scanpy,virshup2024anndata} & 1.11.4 / 0.12.2 & Downstream QC and analysis \\
Scrublet~\cite{wolock2019scrublet} & 0.2.3 & Doublet detection \\
Biopython~\cite{cock2009biopython} & 1.87 & Transcript length / GC from cDNA FASTA \\
statsmodels~\cite{seabold2010statsmodels} & 0.14.5 & Per-gene OLS / Cook's distance \\
edgePython~\cite{robinson2010edger,robinson2010tmm,chen2016edger} & 0.2.5 & Pseudobulk differential composition (edgeR workflow) \\
gseapy (Enrichr)~\cite{fang2023gseapy,kuleshov2016enrichr} & 1.2.1 & Functional enrichment \\
Python & 3.12.3 & --- \\
\hline
\end{tabular}
\end{table}
